\section{方式四：前车轨迹与队列跟随}\label{sec:follow4}
\subsection{方案定位与源码依据}
方式四将队列中的前一台 TurtleBot2 作为目标，复用地图跟随器的位姿输入模式。它与方式二的关系在于：导航和安全门控共享实现，而目标由 TonyPi 背贴标记变为前车在共同地图中的定位位姿。此处按仓库方式四 README 和真实调用链解释实现，不将共享代码包装成另一个独立开发的控制器。2026年10月17日复核主分支提交为\code{43a7771}；前车位姿转发、地图跟随器及轨迹核心与原分析基线一致\cite{repo-current,follow4}。

仓库示例指定后车命名空间\code{tb4}，通过\code{leader_pose_topic:=/tb3/amcl_pose}订阅前车，设置\code{target_source:=pose}与\code{follow_distance:=0.9}。这说明一个明确可复核的\code{tb3}到\code{tb4}配置关系，不代表所提供视频中的车辆已被识别为这些命名空间。用户无法确认画面车辆与方式编号，故视频仅按空间位置命名，不把近端跟随车自动归入方式四。

\subsection{目标身份、坐标与更新节拍}
输入消息为\code{PoseWithCovarianceStamped}，经\code{pose_relay.py}转成\code{PoseStamped}并发布到本车的\code{target_pose}。Relay 检查地图帧、时间戳及协方差，保持原始采样时间，不以转发时刻刷新旧观测。前后车必须共享地图原点；独立里程计\code{odom}不能直接当作共同\code{map}位姿。机器人名称由配置确定，不能从图像中的前后位置推导 ROS 话题身份。

Relay 采用消息回调，没有固定频率重采样器，输出节拍由有效的前车 AMCL 消息决定。地图跟随器使用0.1\,s定时器执行决策，即名义10\,Hz；导航目标则由1\,s时间门限与0.25\,m位移门限共同控制。输入频率、控制决策频率和目标重发频率是三个不同量，不能把10\,Hz写成实际 AMCL 发布频率。
\begin{table}[H]\centering\small
\caption{方式四地图位姿模式的关键接口与默认边界}
\begin{tabularx}{\linewidth}{p{35mm}X}\toprule
项目 & 源码或配置对应关系\\\midrule
示例前车与后车 & 前车\code{/tb3/amcl_pose}；后车命名空间\code{tb4}\\
目标来源 & \code{target_source:=pose}；位姿转发至\code{target_pose}\\
前车观测时效 & 有效年龄为0至0.8\,s；保留原始时间戳\\
定位协方差筛选 & $x/y$方差不超过0.25\,m$^2$，偏航方差不超过0.35\,rad$^2$\\
名义决策周期 & 0.1\,s；未实测运行抖动与消息实际频率\\
路径间隔与近距保护 & 回退间距0.9\,m；欧氏近距保护0.65\,m\\
轨迹采样与跳变 & 最小点间距0.08\,m；大于0.8\,m的跳变拒绝\\\bottomrule
\end{tabularx}\end{table}
表中方差门限是输入筛选条件，不等于实测定位误差。路径回退间距定义在历史折线上，近距保护则使用机器人与前车的欧氏距离；转弯处两者不同，因此必须分别检查，不能用单一“固定车距”概括全部控制行为。

\subsection{历史轨迹回退与串行传播}
轨迹核心\code{Trail}依次保存前车有效地图点，累计长度默认限制为20\,m，并从最新点向后累加弧长，在线段内插值得到跟随目标。设有序点列为$p_0,\ldots,p_n$，期望回退距离为$g$，从末端累计至首次跨越$g$的线段，再求插值系数。该构造沿前车已经测得的路线回退，而非沿前车当前朝向简单偏移。
\begin{equation}
\ell_i=\|p_{i+1}-p_i\|_2,\quad
\alpha=\frac{g-\sum_{j=i+1}^{n-1}\ell_j}{\ell_i},\quad
q=(1-\alpha)p_{i+1}+\alpha p_i,\qquad 0\le\alpha\le1.
\end{equation}
历史长度不足时返回空目标，跟随器等待前车积累路径，不要求后车穿过前车追到路径前方。输出的目标还要经过地图足迹检查，由\code{move_base}/DWA处理导航。该算法保持的是路径几何关系，不构成对所有车辆同时保持刚性队形的全局优化。

若串联$M$条跟随边，每条边都引入观测、决策、导航和执行延迟。在沿路径单向运动的简化条件下，可用$\tau_{\rm tail}\approx\sum_{i=1}^{M}\tau_i$说明末车事件传播延迟的来源；该式不是本次实测延迟，也未证明队列具有串稳定性。定位偏差、转角可见性及目标更新阈值均可能使误差沿队列传播，因此队列长度影响需要同一场地、多次试验的数据，不能仅凭一段视频量化。

\subsection{失联、导航失败与终点处理}
前车位姿过时或协方差不合格时，Relay 不再转发新目标；跟随器在目标时间超过0.8\,s后取消导航并撤销许可。自身定位、TF、代价地图和任务心跳也各有独立守卫。大幅位置跳变被拒绝时不会立即采纳异常新点；旧目标随后可能因时效失效而触发停车。下游速度门控继续检查命令、许可和扫描时效，防止仅停止更新目标而旧速度仍被执行。

导航状态 ABORTED、REJECTED 或 LOST 会锁定失败，等待排故或重置。任务状态为\code{done}且后车接近回退目标至默认0.2\,m以内时，地图跟随器进入完成状态并停止。这是共享地图跟随器的简化完成条件，并非对末车已进入实际终点区域、持续停稳和整队完成的视觉认证。

\subsection{本次视频所支持的过程事实}
视频提供了多车现场布置、TonyPi 抱物和近端底盘处于同一通道中的证据，说明异构平台确实在同一场地参与运行。由于车辆话题身份无法对应，且后续车辆多次被围挡或视角遮挡，本次不填入方式四的实测跟随误差、丢失次数或末车到达率。第\ref{sec:video-evidence}节给出带时间戳的原始截图与证据边界。

要进一步形成方式四的定量结果，所需补充材料是实际启动参数、车辆与命名空间对应表，以及含前后车位姿、任务状态和安全速度的同步日志。视频可核验物理现象，日志则用于确认控制源和计算指标；两类材料共同才能将某次现场行为归因到指定方式。
\source{当前提交\code{43a7771}：\code{turtlebot/follow4/README.md}；\code{follow_map.launch}；\code{pose_relay.py}；\code{map_follower.py}；\code{competition_core/planning.py}。}
